%% ============================================================
%%  §3  The Execution Monad
%% ============================================================

\section{The Execution Monad}
\label{sec:monad}

The semantic foundation of lean-pl-verify is a state-error monad that captures
the observable behaviour of safe Rust and TypeScript programs: they compute
with local state, and they may panic. This section presents the monad
definition, its basic properties, and the forward-backward model for mutable
references.

\subsection{Panic Reasons and Integer Bounds}
\label{sec:types}

Rust programs can panic at runtime for several distinct reasons. We enumerate
them as an inductive type:

\begin{lstlisting}[caption={PanicReason inductive type.},label={lst:panic}]
inductive PanicReason where
  | arithmeticOverflow
  | divisionByZero
  | explicitPanic
  | indexOutOfBounds
\end{lstlisting}

This enumeration is not exhaustive with respect to production Rust (stack
overflow and foreign panics are excluded), but it covers every panic that
can arise in the safe Rust fragment we target.

Integer arithmetic is modelled over \texttt{Int} (Lean's unbounded integer
type), with explicit bounds checks at the point of each operation. The bounds
for 32-bit signed integers are:

\begin{lstlisting}[caption={Integer bounds.},label={lst:bounds}]
namespace IntBounds
  def minI32 : Int := -(2^31)     -- -2147483648
  def maxI32 : Int :=  2^31 - 1  --  2147483647
end IntBounds
\end{lstlisting}

Using \texttt{Int} rather than a fixed-width type (\eg Lean's \texttt{Int32})
simplifies the proof theory significantly: the \texttt{omega} and \texttt{linarith}
tactics operate natively on \texttt{Int}, whereas proofs about \texttt{Int32}
would require constant casts and bitwidth reasoning. The trade-off is that
overflow is a proof obligation rather than a type invariant; every arithmetic
theorem must carry explicit preconditions of the form
$\texttt{minI32} \le v \wedge v \le \texttt{maxI32}$.

\subsection{The RustM Monad}

\begin{lstlisting}[caption={RustM monad definition.},label={lst:rustm}]
def RustM (s : Type) (a : Type) :=
  StateT s (Except PanicReason) a
\end{lstlisting}

\texttt{RustM}~$\sigma$~$\alpha$ is a computation that: reads and writes a
state of type $\sigma$ (modelling local variables), may return a value of
type $\alpha$, and may instead panic with a \texttt{PanicReason}. The
\texttt{StateT} transformer and \texttt{Except} type are both defined in
Lean~4's core library as \texttt{def}s, so they are kernel-transparent.

Monad operations are inherited from the transformer stack:

\begin{lstlisting}[caption={Derived monad operations.},label={lst:monadops}]
-- Fail with a panic reason
def rpanic (r : PanicReason) : RustM s a :=
  fun _ => .error r

-- Safe division: panics on zero denominator
def rdiv (a b : Int) : RustM s Int :=
  if b = 0 then rpanic .divisionByZero
  else pure (a / b)

-- Checked addition: panics on overflow
def raddChecked (a b : Int) : RustM s Int :=
  let v := a + b
  if minI32 <= v /\ v <= maxI32 then pure v
  else rpanic .arithmeticOverflow
\end{lstlisting}

Two monad laws are proved as standalone lemmas and used in the
\ProgramSpec reasoning in \S\ref{sec:programspec}:

\begin{lstlisting}[caption={Monad laws for RustM.},label={lst:monadlaws}]
theorem bind_ok {s a b : Type}
    (f : RustM s a) (g : a -> RustM s b)
    (s : s) (a : a) (s' : s)
    (h : f s = .ok (a, s')) :
    (f >>= g) s = g a s' := by
  simp [bind, StateT.bind, h]

theorem pure_ok {s a : Type} (a : a) (s : s) :
    (pure a : RustM s a) s = .ok (a, s) := rfl
\end{lstlisting}

These look trivial, but their role is structural: they let the \ProgramSpec
combinators in \S\ref{sec:programspec} reason about bound computations without
re-expanding monad definitions in every proof.

\subsection{Kernel Transparency}

A critical design constraint is that \RustM must be \emph{kernel-transparent}:
the Lean kernel must be able to reduce expressions of type
\texttt{RustM\;$\sigma$\;$\alpha$} to their normal forms during type-checking.

This is satisfied by definition: \texttt{StateT}, \texttt{Except}, and all
derived operations are \texttt{def}s with no \texttt{@[extern]} or
\texttt{@[opaque]} attributes. Lean's kernel reduction will unfold all of
them when checking a proof term. The practical consequence is that proofs of
the form \texttt{evalFun [] f 30 args at init |= P} can often be closed
by \texttt{rfl} when \texttt{P} is a \texttt{pureOutput} spec and the
evaluation terminates concretely.

This transparency extends to our \emph{Value} type:

\begin{lstlisting}[caption={Value and Locals types.},label={lst:value}]
inductive Value where
  | int   : Int  -> Value
  | bool_ : Bool -> Value
  | unit  : Value

abbrev Locals := List Value
\end{lstlisting}

Using \texttt{List Value} for local variable storage means that all variable
accesses (\texttt{List.getElem?}, \texttt{List.set}) are also kernel-transparent
standard library functions.

\subsection{The Forward-Backward Model for Mutable References}
\label{sec:fbpair}

Safe Rust frequently passes mutable references as function arguments. Consider
the canonical example from the Aeneas literature~\cite{ho2022aeneas}:

\begin{verbatim}
fn choose<'a, T>(b: bool, x: &'a mut T, y: &'a mut T)
    -> &'a mut T {
  if b { x } else { y }
}
\end{verbatim}

The function returns a mutable reference to one of its inputs. The caller
may then mutate through the returned reference; the mutation must propagate
back to \texttt{x} or \texttt{y} depending on which branch was taken.

In a purely functional model, this is encoded using a pair of functions:

\begin{lstlisting}[caption={FBPair: forward-backward representation.},label={lst:fbpair}]
structure FBPair (a b : Type) where
  fwd : a -> b
  bwd : a -> b -> a

-- forward: b=true => fwd selects x; b=false => fwd selects y
-- backward: propagates updated value back to the chosen input
\end{lstlisting}

The \texttt{fwd} function (forward) computes the result of selecting a branch.
The \texttt{bwd} function (backward) takes the original input and the updated
output and returns the updated input — modelling the lifetime-end propagation.

Four identity laws are proved:

\begin{lstlisting}[caption={FBPair identity laws.},label={lst:fblaws}]
theorem fwd_true (x y : a) :
    choose_fwd true x y = x := rfl

theorem fwd_false (x y : a) :
    choose_fwd false x y = y := rfl

theorem identity_true (x y x' : a) :
    choose_bwd true x y x' = x' := rfl

theorem identity_false (x y y' : a) :
    choose_bwd false x y y' = y' := rfl
\end{lstlisting}

These four theorems constitute the \emph{forward-backward identity laws}:
applying the backward function with the forward result returns the input
unchanged. They are the functional analogue of the borrow-lifetime protocol
and suffice to reason about any function that selects a borrow based on
a pure condition.

\subsection{Comparison with Dijkstra Monads}

The LOOM~\cite{loom2026} and F$^\star$~\cite{fstar2016} approaches use Dijkstra
monads~\cite{dijkstraMonads2019}, which are indexed by pre- and post-conditions
in a way that generates verification conditions automatically. \RustM is
deliberately simpler: it is a plain monad without WP-index.

The choice reflects our use case. Dijkstra monads excel when the goal is to
automatically generate proof obligations from program code. Our goal is
different: we want to state explicit, human-readable specifications
(\ProgramSpec) and prove them by hand using Lean tactics. A plain monad
leaves the specification structure fully in the user's control and avoids
the complexity of a WP-monad universe hierarchy. Whether Dijkstra monads
would further automate the loop invariant proofs in \S\ref{sec:rust-proofs}
is an interesting open question we leave for future work.
